% Gerado de docs/monografia/curadoria.md por scripts/monografia_latex.py.
% Não edite aqui: edite o .md e gere de novo.

\chapter{\texorpdfstring{\mbox{CURADORIA} \mbox{DO} \mbox{CATÁLOGO} \mbox{DE} \mbox{CONCEPÇÕES} \mbox{ERRÔNEAS}}{CURADORIA DO CATÁLOGO DE CONCEPÇÕES ERRÔNEAS}}\label{cap:4}

\section{Problema da procedência}

O sistema descrito nos capítulos anteriores formula, a cada atividade respondida, uma hipótese com consequência imediata: a de que a pessoa provavelmente sustenta determinada concepção errônea. É essa hipótese que seleciona o percurso de refutação a ser apresentado.

A validade da hipótese depende de uma questão anterior, que o software não responde: a origem de cada concepção do catálogo. Se uma concepção tiver sido incluída apenas por parecer plausível ao autor, o sistema não estará diagnosticando uma concepção errônea documentada, e sim produzindo dados sobre uma suposição, com a mesma aparência de rigor do restante. Por essa razão, o projeto adota a seguinte regra: uma concepção sem procedência é tratada como hipótese do autor e não é incluída no catálogo como achado.

Este capítulo trata dessa questão, que pode ser verificada sem a participação de estudantes. O instrumento não foi avaliado com estudantes, limitação declarada no Capítulo 7. O que se sustenta aqui é mais restrito e inteiramente verificável: o sistema afirma sobre concepções errôneas apenas o que a literatura registra e identifica explicitamente os casos em que não há registro.

\section{Graus de procedência}

Cada concepção do catálogo possui um campo de procedência (\texttt{provenance}), que assume um de cinco valores. O grau não descreve a importância da concepção nem a confiança do autor nela: descreve o tipo de obra que a sustenta (Quadro~\ref{qua:4-1}).

\begin{quadro}[htbp]
\caption{Graus de procedência}\label{qua:4-1}
\footnotesize
\begin{tabularx}{\textwidth}{|l|X|}
\hline
\textbf{grau} & \textbf{significado} \\ \hline
\texttt{D} & documentada: achado empírico, com fonte e posição localizável \\ \hline
\texttt{L} & registrada como erro comum em obra técnica \\ \hline
\texttt{T} & transposta de outra linguagem, com a analogia declarada \\ \hline
\texttt{N} & o erro consta de norma, como o padrão de codificação CERT C \cite{sei2025} e a CWE, \textit{Common Weakness Enumeration} \cite{mitre2026}, mas a concepção não \\ \hline
\texttt{E} & elaborada a partir do modelo de máquina: hipótese a testar \\ \hline
\end{tabularx}
\fonte{Elaborado pelo autor.}
\end{quadro}

Duas propriedades desse campo são relevantes para o método.

A primeira é que o grau constitui uma afirmação sobre a literatura, e não sobre o uso da concepção no sistema. Uma mesma concepção pode aparecer em mais de um tópico do domínio (o que ocorre com quatorze delas), e o texto que a contextualiza varia de um tópico para outro; o grau, porém, deve ser o mesmo. Se uma concepção tem grau \texttt{D} em um tópico e \texttt{L} em outro, uma das duas afirmações é falsa, e essa diferença corresponde, em um trabalho acadêmico, à diferença entre citar e alegar. A auditoria encontrou três casos desse tipo (Seção 4.6.5), e um teste automatizado impede a divergência (Quadro~\ref{qua:4-7}).

A segunda é que o grau \texttt{E} é legítimo, desde que declarado. Uma hipótese do autor pode integrar o catálogo, mas não pode ser apresentada como achado. As três hipóteses do catálogo são identificadas como tais no próprio campo de fonte, e dois testes independentes asseguram que nenhuma delas apareça nos tópicos selecionados para a coleta de dados.

\section{Classes de fonte}

O catálogo apoia-se em dois tipos de obra. A diferença entre eles não é de qualidade, mas do tipo de afirmação que cada um sustenta. Tratá-los como equivalentes constituiria um defeito metodológico; a declaração da classe, concepção por concepção, é o que torna o catálogo auditável.

\subsection{Classe 1: literatura de concepções errôneas (23 concepções)}

A primeira classe compreende estudos empíricos com estudantes e revisões que os agregam. A obra principal é o Apêndice A de \citeonline{sorva2012}, cuja Tabela A.1 relaciona 162 concepções errôneas levantadas na literatura de ensino de programação introdutória, cada uma com tópico, descrição e fonte. A ela se somam os estudos que a própria tabela atribui: du Boulay (\citeyear{duboulay1986}), \citeonline{soloway1982}, \citeonline{putnam1986}, \citeonline{sleeman1986}, \citeonline{samurcay1989}, \citeonline{holland1997}, \citeonline{ma2007}, \citeonline{doukakis2007} e \citeonline{kaczmarczyk2010}. Uma obra externa à Tabela A.1 também sustenta uma das concepções: \citeonline{hermans2018} registram a ideia de que a variável armazena vários valores ou ``lembra'' os anteriores e a investigam em um experimento com 496 iniciantes.

Essa classe sustenta que a concepção foi observada em estudantes por pesquisadores que a investigaram.

\subsection{Classe 2: corpus de praticante (34 concepções)}

A segunda classe é constituída pelo FAQ (\textit{Frequently Asked Questions}, perguntas frequentes) do grupo comp.lang.c, compilado em 1990 a partir do grupo de Usenet de mesmo nome, aprimorado de forma episódica desde então e publicado em livro por \citeonline{summit1996}.

A inclusão dessa classe requer justificar por que um FAQ de linguagem constitui evidência sobre concepções errôneas, e não apenas sobre a linguagem C. A justificativa está no critério de inclusão adotado, que pode ser conferido por qualquer leitor: é incluída somente a pergunta cujo título enuncia a concepção na voz de quem a sustenta, em primeira pessoa (Quadro~\ref{qua:4-2}).

\begin{quadro}[htbp]
\caption{Exemplos de perguntas do FAQ cujo título enuncia a concepção na voz de quem a sustenta}\label{qua:4-2}
\footnotesize
\begin{tabularx}{\textwidth}{|l|X|}
\hline
\textbf{pergunta} & \textbf{título} \\ \hline
6.2 & ``But I heard that \texttt{char a[]} was identical to \texttt{char *a}.'' \\ \hline
6.9 & ``Someone explained to me that arrays were really just constant pointers.'' \\ \hline
6.4b & ``So arrays are passed by reference, even though the rest of C uses pass by value?'' \\ \hline
4.11 & ``Does C even have ``pass by reference''?'' \\ \hline
\end{tabularx}
\fonte{Elaborado pelo autor, com títulos de perguntas de \citeonline{summit1996}.}
\end{quadro}

O catálogo, portanto, não infere a concepção a partir de uma explicação técnica: registra a formulação de quem a sustenta. Para conferir a inclusão, basta consultar a pergunta e avaliar se o título enuncia a concepção ou apenas o assunto.

Do ponto de vista metodológico, o FAQ é tratado como fonte documental: um corpus de confusões autorrelatadas, de ocorrência natural, selecionadas por recorrência ao longo de mais de três décadas. O tratamento dado a esse corpus é o da análise documental.

Essa classe sustenta que a confusão ocorre e é recorrente entre aprendizes da linguagem C. Não sustenta afirmações de prevalência: a população é autosselecionada (pessoas que escrevem para um grupo de Usenet), e a frequência das perguntas reflete o julgamento do editor, e não uma medição. Por isso, nenhuma afirmação de proporção é extraída dessa fonte.

\subsection{Necessidade da segunda classe}

Seria razoável exigir que todo o catálogo se apoiasse na classe 1, de natureza empírica, e que a classe 2 fosse dispensada. Essa exigência foi testada contra a Tabela A.1 de \citeonline{sorva2012}, item por item, e verificou-se que não pode ser satisfeita.

A Tabela A.1 possui um tópico dedicado a referências (\texttt{Refs}, com 17 dos 162 itens) e itens sobre memória e tempo de vida distribuídos pelos tópicos \texttt{VarAssign}, \texttt{ObjState} e \texttt{Misc}. Faltam-lhe, contudo, os itens relevantes para este trabalho, como mostra a contagem do Quadro~\ref{qua:4-3}, reprodutível por busca de termos no texto da tabela.

\begin{quadro}[htbp]
\caption{Ocorrências de termos relativos à memória na Tabela A.1}\label{qua:4-3}
\footnotesize
\begin{tabularx}{\textwidth}{|X|X|}
\hline
\textbf{termo buscado na Tabela A.1} & \textbf{ocorrências} \\ \hline
\texttt{malloc}, \texttt{free}, \texttt{deallocat}, \texttt{garbage} & 0 \\ \hline
\texttt{heap}, \texttt{stack} & 0 \\ \hline
\texttt{arithmetic}, \texttt{decay}, \texttt{dangling}, \texttt{sizeof} & 0 \\ \hline
\texttt{null} & 1: o item 54, ``Null model of recursion'' \cite[p.~362]{sorva2012}, que não trata de ponteiro nulo \\ \hline
\texttt{lifetime} & 1: o item 5, ``Difficulties understanding the lifetime of values'' \cite[p.~359]{sorva2012} \\ \hline
\texttt{memory} & 7 \\ \hline
\end{tabularx}
\fonte{Elaborado pelo autor, por busca no texto de \citeonline{sorva2012}.}
\end{quadro}

Não há, portanto, nenhum item sobre \texttt{NULL}, sobre \texttt{malloc} e \texttt{free}, sobre aritmética de ponteiros, sobre decaimento de array ou sobre a distinção entre pilha e monte.

Essa ausência não decorre de descuido da obra: o catálogo de \citeonline{sorva2012} refere-se à programação introdutória em Java, Python e BASIC, linguagens em que esses construtos não existem. A literatura de concepções errôneas produzida nessas linguagens não poderia documentar confusões relativas justamente ao que distingue a linguagem C.

Conclui-se, assim, não que as 34 concepções da classe 2 sejam frágeis, mas que um catálogo de concepções errôneas para C requer uma segunda classe de fonte, a qual deve ser declarada como tal, com a indicação do que sustenta e do que não sustenta. Essa declaração é feita, concepção por concepção, na coluna de classe do Apêndice B.

\section{Procedimento de conferência}

O padrão adotado pelo projeto, registrado no cabeçalho do primeiro catálogo de tópico, é o de que as fontes sejam conferidas na obra primária, e não citadas de memória. O resultado da aplicação desse padrão a cada concepção é apresentado no Quadro~\ref{qua:4-4}.

O quadro é exaustivo e mutuamente exclusivo: cada uma das 61 concepções figura em exatamente uma linha, que registra a verificação mais forte obtida para ela. A soma é, por construção, igual a 61. Um quadro que contasse as concepções que citam cada obra teria soma maior, pois uma concepção pode citar mais de uma obra, e não permitiria determinar quantas ficaram sem conferência.

\begin{quadro}[htbp]
\caption{Forma de conferência de cada uma das 61 concepções}\label{qua:4-4}
\footnotesize
\begin{tabularx}{\textwidth}{|X|l|}
\hline
\textbf{forma de conferência} & \textbf{concepções} \\ \hline
conferida no texto integral de \citeonline{sorva2012}, obtido do repositório aberto da Universidade Aalto (handle \texttt{123456789/3534}, bundle TEXT); a frase em inglês citada entre aspas foi localizada no texto da obra & 21 \\ \hline
conferida no índice do FAQ comp.lang.c: as seções 1, 4, 5, 6 e 7 foram consultadas em \texttt{c-{}faq.com}; as perguntas citadas existem, e os títulos coincidem com os citados no catálogo & 32 \\ \hline
citação abreviada no catálogo, com a forma integral localizada na obra & 1 \\ \hline
cita \citeonline{sorva2012} com página e item, mas sem frase entre aspas, o que impede a verificação pelo procedimento automatizado & 1 \\ \hline
conferência parcial: a seção 11 do FAQ não foi incluída na curadoria & 1 \\ \hline
não conferida na curadoria: cita apenas normas (CERT C e CWE), conferidas à parte (Seção 4.7) & 1 \\ \hline
sem item correspondente na literatura, com a busca declarada em campo próprio & 2 \\ \hline
sem fonte, por definição: hipótese declarada como tal & 2 \\ \hline
total & 61 \\ \hline
\end{tabularx}
\fonte{Elaborado pelo autor.}
\end{quadro}

As duas primeiras linhas somam as 53 concepções conferidas na obra primária. As seis últimas correspondem ao que a curadoria não afirma ter conferido, com a respectiva justificativa.

Os estudos agregados pelo Apêndice A de \citeonline{sorva2012}, a saber, du Boulay (\citeyear{duboulay1986}), \citeonline{samurcay1989}, \citeonline{sleeman1986}, \citeonline{soloway1982}, \citeonline{putnam1986}, \citeonline{doukakis2007}, \citeonline{ma2007} e \citeonline{kaczmarczyk2010}, não figuram em linha própria porque não foram consultados diretamente. As cinco concepções que os citam fazem-no \textit{apud} \citeonline{sorva2012}, com página e item, e estão contadas na primeira linha do quadro. A Seção 4.7 detalha esse ponto.

A conferência é automatizada sempre que possível: o programa que gera o Apêndice B extrai a frase entre aspas de cada citação e a procura no texto da obra, de modo que o resultado não depende do que está declarado no campo de fonte. Quando o programa é executado sem o texto da obra, o número de concepções conferidas cai de 53 para 32, e o apêndice passa a indicar, linha por linha, que a conferência não foi realizada.

\section{Resultado da curadoria}

O Quadro~\ref{qua:4-5} sintetiza o resultado da curadoria.

\begin{quadro}[htbp]
\caption{Resultado da curadoria}\label{qua:4-5}
\footnotesize
\begin{tabularx}{\textwidth}{|X|X|}
\hline
\textbf{indicador} & \textbf{valor} \\ \hline
concepções distintas no catálogo & 61 \\ \hline
conferidas na obra primária & 53 \\ \hline
graus & \texttt{D} 6, \texttt{L} 33, \texttt{T} 16, \texttt{N} 3, \texttt{E} 3 \\ \hline
hipóteses \texttt{E} & 3, todas fora dos tópicos selecionados para a coleta \\ \hline
\end{tabularx}
\fonte{Elaborado pelo autor.}
\end{quadro}

\section{Achados da auditoria}

Os achados descritos nesta seção têm interesse que ultrapassa o catálogo, pois dizem respeito ao método de construção de um catálogo a partir da literatura, e não à linguagem C.

\subsection{Buscas realizadas com vocabulário ou item inadequado}

Em três concepções, a busca inicial na obra não localizou o item correspondente, embora ele existisse (Quadro~\ref{qua:4-6}).

\begin{quadro}[htbp]
\caption{Buscas que não localizaram o item e o item existente na obra}\label{qua:4-6}
\footnotesize
\begin{tabularx}{\textwidth}{|X|X|X|}
\hline
\textbf{concepção} & \textbf{registro anterior no catálogo} & \textbf{item existente na obra} \\ \hline
\texttt{pointer\_\allowbreak{}has\_\allowbreak{}no\_\allowbreak{}address} & ``procurei no catálogo do Sorva e não está lá: nenhum dos 162 itens afirma que a referência não ocupa espaço'' & a segunda frase do item 14: ``It is not stored inside the computer'' \cite[p.~360]{sorva2012}. A busca utilizou os termos ``referência'' e ``espaço'', ao passo que o item trata de variável \\ \hline
\texttt{array\_\allowbreak{}confused\_\allowbreak{}with\_\allowbreak{}cell} & ``não tem fonte empírica, entra como hipótese'' & o item 151, ``Confusion between an array and its cell'' \cite[p.~367]{sorva2012}, e o item 9, ``A variable can hold multiple values at a time'' \cite[p.~359]{sorva2012}. A busca utilizou o termo ``array'', ao passo que o item 9 trata de variável e o item 151 pertence ao tópico \texttt{Misc} \\ \hline
\texttt{assignment\_\allowbreak{}symmetric} & citava o item 11, com a aproximação declarada: ``o 11 é sobre a direção invertida, e esta concepção é sobre os dois lados'' & o item 12, imediatamente posterior ao 11: ``Primitive assignment works both directions (swaps)'' \cite[p.~359]{sorva2012}. O item corresponde ao rótulo da concepção \\ \hline
\end{tabularx}
\fonte{Elaborado pelo autor, com trechos de \citeonline{sorva2012}.}
\end{quadro}

Nos três casos, a busca refeita com outro vocabulário localizou o item, e é esse o item citado no catálogo, com o grau que ele sustenta: \texttt{T} para a primeira concepção e \texttt{D} para as outras duas. Nenhuma das três classificações depende de julgamento intuitivo.

Desses casos decorre uma recomendação de método, que pode ser replicada: a busca deve ser orientada pelo fenômeno, e não pelo termo da linguagem em estudo. Em um catálogo de concepções errôneas da linguagem X, o item procurado pode estar redigido no vocabulário da linguagem Y, o que necessariamente ocorre quando o catálogo foi produzido em Y.

O terceiro caso acrescenta um segundo princípio, relativo à transparência da citação: foi a aproximação declarada que permitiu localizar o erro. A citação indicava em que o item diferia da concepção, e essa diferença apontou para o item vizinho. Uma citação que omite a discrepância não oferece esse indício e permanece incorreta sem que se perceba.

\subsection{Citação associada a item incorreto}

A citação da concepção \texttt{pointer\_\allowbreak{}is\_\allowbreak{}data} remetia ao item 14 de \citeonline[p.~360]{sorva2012}, que trata de a variável ser um par de nome e valor, e não de ponteiros. A concepção consta do item 65: ``A (non-primitive Java) variable does not hold a reference but a set of object properties'' \cite[p.~362]{sorva2012}. A obra citada estava correta, mas o item não, e esse tipo de erro somente pode ser identificado com a consulta à obra.

O mesmo tipo de erro ocorre na identificação das obras mediadas. A Tabela A.1 atribui quatorze dos itens utilizados a ``\citeonline{ma2007}'', e um artigo publicado no simpósio SIGCSE no mesmo ano e sobre o mesmo tema, assinado por Ma e outros três autores, pode ser facilmente confundido com essa obra; a bibliografia de \citeonline{sorva2012}, entretanto, identifica ``\citeonline{ma2007}'' como uma tese de doutorado de autoria única, de outra instituição. A citação de uma obra mediada exige, portanto, a consulta à bibliografia da obra mediadora, não bastando o reconhecimento do sobrenome e do ano.

\subsection{Fontes que sustentam mais de uma concepção}

As concepções \texttt{memory\_\allowbreak{}is\_\allowbreak{}by\_\allowbreak{}name} e \texttt{pointer\_\allowbreak{}has\_\allowbreak{}no\_\allowbreak{}address} derivam das duas frases do mesmo item 14 de \citeonline[p.~360]{sorva2012}. No FAQ de \citeonline{summit1996}, a pergunta 7.21 sustenta três concepções de um tópico; a 7.22, duas de outro; e a 4.9, duas de um terceiro.

Não se trata de economia de fontes: uma resposta de FAQ frequentemente trata de mais de uma confusão, e separá-las constitui trabalho de curadoria, e não apenas de leitura.

\subsection{Citações não sustentadas pela obra citada}

A regra aplicada é a seguinte: se a fonte não pode ser verificada, procura-se o item correspondente na obra; se ele não existir, a citação não permanece no catálogo. Cinco casos foram identificados:

\begin{itemize}
\item Atribuição incorreta. A citação de \texttt{declaration\_\allowbreak{}allocates} atribuía a concepção a \citeonline{kaczmarczyk2010}, com página e código de tabela que somente poderiam ser conferidos com a consulta ao artigo. O item que registra essa concepção, o de número 92, atribui-a a \citeonline{ma2007}, e é essa a atribuição citada no catálogo, \textit{apud} \citeonline[p.~364]{sorva2012}. Como a concepção pertence a um dos tópicos selecionados para a coleta de dados, a atribuição tem consequência prática.
\item Concepção sem item na obra. Para \texttt{memory\_\allowbreak{}is\_\allowbreak{}uniform}, a citação indicava página e tabela de um artigo sem cópia de acesso aberto. Na busca item por item no Apêndice A, não se encontrou item sobre a memória como região única, sem regras distintas de tempo de vida; o mais próximo, o item 5 \cite[p.~359]{sorva2012}, refere-se à \textit{dificuldade} com o tempo de vida, que é outro fenômeno. No catálogo, a concepção é uma hipótese declarada, de grau \texttt{E}.
\item Refutação no lugar da fonte. Para \texttt{declaration\_\allowbreak{}zeroes}, a fonte citada era \citeonline{kernighan1988}, obra que afirma ser falsa a concepção; trata-se, portanto, da refutação, e não do registro de que alguém a sustenta. O catálogo distingue os dois papéis: a fonte é o item 10 \cite[p.~359]{sorva2012}, e a obra de \citeonline{kernighan1988} é registrada como refutação.
\item Código interno de artigo não consultado. A citação de \texttt{no\_\allowbreak{}value\_\allowbreak{}no\_\allowbreak{}memory} continha o código ``(PVR2)'', pelo qual \citeonline{kaczmarczyk2010} identificam o achado e que somente pode ser verificado com a consulta ao artigo. O catálogo cita o item 22, \textit{apud} \citeonline[p.~360]{sorva2012}.
\item Aproximação declarada que indicou o item correto. É o caso de \texttt{assignment\_\allowbreak{}symmetric}, descrito na Seção 4.6.1 por constituir também um resultado de método.
\end{itemize}

\subsection{Procedência divergente entre tópicos}

Quatorze concepções aparecem em mais de um tópico, e cada ocorrência possui campo de fonte próprio, preenchido manualmente. Três delas divergiam quanto ao grau e, em dois desses três casos, uma das ocorrências estava correta: citava página e item, ao passo que a ocorrência do outro tópico citava o estudo original sem declarar a mediação. A conferência havia sido realizada em um tópico e não fora propagada ao outro, sem que o sistema acusasse a inconsistência.

O achado relevante não se restringe a essas três concepções: um catálogo mantido manualmente em arquivos paralelos tende a divergir sem que isso seja percebido. As ocorrências estão uniformizadas, e um teste automatizado impede novas divergências (Quadro~\ref{qua:4-7}).

\section{Integridade das citações}

Esta seção explicita o que foi lido no original, em contrapartida ao fato de o catálogo se apoiar em uma obra que, por sua vez, agrega dezenas de outras.

Foram lidos no original o texto integral de \citeonline{sorva2012}, obtido do repositório institucional aberto da Universidade Aalto; as perguntas citadas do FAQ do comp.lang.c, consultadas em \texttt{c-{}faq.com}; e o artigo de \citeonline{hermans2018}. As citações dessas obras são diretas, e a frase citada foi localizada no texto.

Os estudos agregados pelo Apêndice A, a saber, du Boulay (\citeyear{duboulay1986}), \citeonline{samurcay1989}, \citeonline{sleeman1986}, \citeonline{soloway1982}, \citeonline{putnam1986}, \citeonline{doukakis2007}, \citeonline{ma2007} e \citeonline{kaczmarczyk2010}, são citados \textit{apud} \citeonline{sorva2012}, com página e item. Os artigos originais não foram consultados, por se encontrarem em acervos de acesso restrito; a verificação realizada limitou-se à existência do registro bibliográfico, sem acesso ao texto. Cinco concepções do catálogo citam esses estudos, e cada uma declara a mediação no próprio campo de fonte.

Em consequência, nenhuma afirmação deste trabalho atribui a esses artigos conteúdo que não conste do item do Apêndice A que os cita.

As normas CERT C \cite{sei2025} e CWE \cite{mitre2026}, que sustentam as três concepções de grau \texttt{N}, foram conferidas separadamente da curadoria, na conferência das referências deste trabalho: as três entradas da CWE (252, 415 e 476) e as quatro regras do CERT C (ERR33-C, EXP34-C, MEM30-C e MEM31-C) correspondem ao que o catálogo cita. Os Quadros~\ref{qua:4-4} e~\ref{qua:4-5} referem-se à conferência realizada na curadoria, na qual a concepção que cita apenas normas ficou sem conferência e a seção 11 do FAQ não foi incluída.

Duas observações de forma delimitam o que se pode afirmar. A primeira é que a versão on-line do FAQ é mais curta do que o livro de \citeonline{summit1996}, como informa o próprio autor na página do FAQ; as citações deste trabalho referem-se à versão on-line, por número de pergunta, que foi a versão consultada. A segunda é que os números de página do Apêndice A referem-se ao texto extraído do PDF disponível no repositório mencionado e foram calibrados com base nos marcadores de página do próprio texto.

\section{Limites da curadoria}

A delimitação do alcance da curadoria integra o método. Cinco limites devem ser considerados:

\begin{enumerate}
\item Não há afirmação de prevalência. Nem a classe 1 nem a classe 2 autorizam afirmar que determinada proporção dos estudantes sustenta uma dada concepção. O catálogo afirma que a concepção existe e é recorrente.
\item Não há afirmação sobre a eficácia do instrumento. O fato de as concepções estarem documentadas não implica que o sistema as detecte adequadamente nem que a refutação produza efeito, o que exigiria coleta de dados com participantes.
\item Três concepções permanecem como hipóteses. São elas \texttt{memory\_\allowbreak{}is\_\allowbreak{}uniform}, \texttt{type\_is\_a\_hint} e \texttt{memoria\_\allowbreak{}e\_\allowbreak{}do\_\allowbreak{}computador}, declaradas como tais e ausentes dos tópicos de coleta. A auditoria indica o procedimento a adotar antes de aceitá-las: refazer a busca com outro vocabulário, procedimento que localizou o item de duas concepções antes consideradas hipóteses (Seção 4.6.1) e que, para essas três, confirmou a inexistência de item.
\item Uma concepção de grau \texttt{T} não cita nenhuma obra. A concepção \texttt{heap\_\allowbreak{}dies\_\allowbreak{}with\_\allowbreak{}function} apoia-se apenas na analogia declarada: em Java e em Python, a memória alocada por uma função e que deixa de ser alcançável é liberada automaticamente. A busca foi refeita com o vocabulário do fenômeno e registrada, sem que se localizasse item correspondente. O próprio campo de fonte indica que a concepção requer uma fonte ou deve ser reclassificada para o grau \texttt{E}.
\item Uma citação de \citeonline{sorva2012} não contém frase entre aspas. Trata-se da citação de \texttt{uninitialized\_\allowbreak{}is\_\allowbreak{}null}, que, embora indique página e item, não pode ser verificada pelo procedimento automatizado descrito na Seção 4.4.
\end{enumerate}

\section{Verificação automatizada do catálogo}

O Apêndice B não foi redigido manualmente: é gerado a partir do catálogo por um comando. A digitação de 61 linhas favoreceria exatamente a divergência entre documento e dado que este capítulo se propõe a controlar, e a Seção 4.6.5 mostra que essa divergência ocorre mesmo quando o risco é conhecido por quem mantém os arquivos.

O catálogo é submetido a verificações automatizadas, que fazem a suíte de testes falhar caso a propriedade correspondente deixe de valer. As principais são apresentadas no Quadro~\ref{qua:4-7}.

\begin{quadro}[htbp]
\caption{Principais verificações automatizadas do catálogo}\label{qua:4-7}
\footnotesize
\begin{tabularx}{\textwidth}{|X|X|}
\hline
\textbf{verificação} & \textbf{situação rejeitada} \\ \hline
grau uniforme entre tópicos & a mesma concepção com grau \texttt{D} em um tópico e \texttt{L} em outro \\ \hline
grau \texttt{D} com posição & grau \texttt{D} sem página nem item na fonte \\ \hline
grau \texttt{E} declarado & hipótese que não se identifica como hipótese \\ \hline
grau \texttt{E} fora dos tópicos de coleta & hipótese em tópico selecionado para a coleta (verificado por dois caminhos independentes) \\ \hline
busca sem resultado não conta como conferência & linha que indique ``conferida'' para concepção que declara, em campo próprio, não haver item na literatura \\ \hline
busca sem resultado declarada em todas as ocorrências & campo que a declara presente em um tópico e ausente em outro \\ \hline
ausência de chave duplicada no YAML & campo sobrescrito por repetição, sem aviso \\ \hline
consistência entre capítulo e catálogo & número afirmado neste capítulo que o catálogo deixou de sustentar \\ \hline
\end{tabularx}
\fonte{Elaborado pelo autor.}
\end{quadro}

A reprodução da curadoria requer um único comando:

\begin{center}\small\ttfamily
python~scripts/gera\_curadoria.py~CAMINHO/DO/sorva2012.txt
\end{center}

O argumento é o texto integral da obra, obtido do repositório da Universidade Aalto. O arquivo não integra o repositório do projeto por se tratar de obra de terceiro; na sua ausência, o programa é executado normalmente, e o apêndice passa a indicar, linha por linha, que a conferência não foi realizada.

A contribuição deste capítulo reside nessa verificabilidade, e não no número de concepções: um catálogo de 61 concepções errôneas sobre ponteiros em C, com duas classes de fonte declaradas, posição localizável sempre que o grau a exige, mediação de citação declarada onde existe, e conferível em seu conjunto por um único comando.
